% Generated draft: requires mathematical and source review.
\documentclass[11pt]{article}
\usepackage[margin=25mm]{geometry}
\usepackage{fontspec}
\setmainfont{Latin Modern Roman}
\usepackage{amsmath,amssymb,mathtools,unicode-math}
\setmathfont{Latin Modern Math}
\usepackage{enumitem,needspace,xurl,hyperref}
\setlength{\parindent}{0pt}
\setlength{\parskip}{5pt}
\setlength{\emergencystretch}{3em}
\newcommand{\N}{\mathbb{N}}
\newcommand{\Z}{\mathbb{Z}}
\newcommand{\Q}{\mathbb{Q}}
\newcommand{\R}{\mathbb{R}}
\newcommand{\C}{\mathbb{C}}
\newcommand{\F}{\mathbb{F}}
\newcommand{\A}{\mathbb{A}}
\newcommand{\PP}{\mathbb{P}}
\newcommand{\E}{\mathbb{E}}
\newcommand{\eps}{\varepsilon}
\newcommand{\dd}{\,\mathrm{d}}
\begin{document}
\section*{UnsolvedMath ID: 30005439}
\textbf{Source problem code:} catalogue entry
\subsection{Definitions and mathematical statement}
Under the usual left-operation convention, a \emph{cycle set} is a set \(X\) with a binary operation \(\cdot:X\times X\to X\) such that each map \(L_x:X\to X\), \(L_x(y)=x\cdot y\), is bijective and
\[
(x\cdot y)\cdot(x\cdot z)=(y\cdot x)\cdot(y\cdot z)
\qquad(x,y,z\in X).
\]
Under this convention, non-degeneracy means that \(x\mapsto x\cdot x\) is bijective. These conventions are stated conditionally: the supplied excerpt does not give the defining axioms or operation conventions. It also does not specify whether \(X\) must be finite or nonempty.

A \emph{commutative semigroup operation} on a set \(A\) is a binary operation \(\oplus:A\times A\to A\) satisfying
\[
(a\oplus b)\oplus c=a\oplus(b\oplus c),
\qquad a\oplus b=b\oplus a
\qquad(a,b,c\in A).
\]

\textbf{Qualified problem statement.} For cycle sets in the intended source class, find further deformations into cycloids and determine the commutative semigroup operations on the associated structure object \(M(X)\) induced by these deformations. The excerpt gives the generator-level expression
\[
x\oplus y=(x\cdot y)x,
\]
but does not define the juxtaposition product, identify which operation \(\cdot\) is used after deformation, or specify how this expression extends to the whole structure object.

The supplied clean statement calls \(M(X)\) the structure \emph{monoid}, whereas the quoted original calls it the structure \emph{group} and previously refers to semigroup structures on \(S(X)\). No presentations or identifications of these objects are supplied. Consequently, the carrier of the requested operations cannot be fixed from this record alone. Likewise, the axioms of a cycloid and an L-algebra, the meaning of deformation, and the criterion for a deformation to count as further or different are missing. These qualifications prevent a fully specified classification problem.

The source states that non-degenerate cycle sets can be deformed into L-algebras, that both types of structure are cycloids, and that they extend to different semigroup structures on \(S(X)\). It cites cabling as an example transforming a non-degenerate cycle set into another non-degenerate cycle set. Whether the requested new deformations are restricted to non-degenerate inputs is not explicit. The surrounding passage anticipates useful normal forms for non-involutive structure invariants and a better understanding of combinatorial properties such as growth; these are stated expectations, not established consequences or separately specified mathematical targets.
\subsection{Short English statement}
Find additional ways to deform cycle sets into cycloids and describe the commutative semigroup operations they induce on the intended structure object. A precise formulation requires the missing deformation and cycloid definitions, together with clarification of the conflicting references to the structure monoid \(M(X)\), structure group \(M(X)\), and \(S(X)\).
\subsection{Sources}
\begin{itemize}
\item Supplied local UnsolvedMath registry record, permanent ID \(30005439\), titled \emph{Cycle Set Deformations into Cycloids}; exact source code \texttt{catalogue entry}. Local registry locations: \texttt{data/unsolvedmath/problems.json} and the supplementary \texttt{catalog/problems.json}. No independent inspection of these files is asserted.
\item \emph{Mini-Workshop: Skew Braces and the Yang--Baxter Equation} (2023), Oberwolfach Reports, pp. 537--563, DOI \texttt{10.4171/owr/2023/9}: \url{https://doi.org/10.4171/owr/2023/9}. The supplied original excerpt identifies Problem 12 and includes the page marker 547. This DOI link is also the supplied literature source labelled ``original OWR report.''
\item The excerpt attributes the cabling technique to Lebed, Ramirez, and Vendramin, citing its reference [4]. The bibliographic entry for [4] is not supplied; no additional citation or URL is inferred. No proposer is identified in the record.
\item Literature assessment dated exactly \texttt{2026-08-22}: the supplied curation labels the problem open and reports that no complete classification or resolution was located. This is supplied dated curation, not an independently verified literature assessment. The record also reports direct verification against the official report and assigns the label \texttt{corrected\_verified}; these are supplied verification claims, not verification performed for this draft.
\end{itemize}
\end{document}
